\section{Spectral properties of compact linear operators on normed spaces}


\begin{exercise}{2}
Let $X, Y$ and $Z$ be normed spaces, and let $T: X \to Y$ and $S: Y \to Z$.
If $T$ and $S$ are compact linear operators, show that $S T: X \to Z$ is a compact linear operator.
\end{exercise}
\begin{proof}
We have that $T$ is bounded (8.1-2) so that for a bounded set $B \subseteq X$, $T(B)$ is bounded, thus $S T(B)$ is relatively compact, meaning that $S T$ is compact (8.3-2), as required.
\end{proof} 

\begin{exercise}{3}
If $T$ is a compact linear operator, show that for any given number $k > 0$ there are at most finitely many linearly independent eigenvectors of $T$ corresponding to eigenvalues of absolute value greater than $k$.
\end{exercise}
\begin{proof}
In 8.3-1 we proved that there are finitely many eigenvalues with $\absoluteValue{\lambda} > k$.
Moreover, in 8.3-3 we proved that the null space of $T_\lambda$ (for $\lambda \neq 0$ is finite dimensional.
Since this null space corresponds to the eigenvectors, the union of all eigenvectors (for a fixed $k$) is finite, completing the proof.
\end{proof} 

\begin{exercise}{4}
Let $T_j: X_j \to X_{j+1}$, $j = 1, 2, 3$ be bounded linear operators on normed spaces.
If $T_2$ is compact, show that $T = T_3 T_2 T_1: X_1 \to X_4$ is compact.
\end{exercise}
\begin{proof}
Lemma 8.3-2 tells us that $T_2 T_1$ is compact.
Applying the same principle to $T_3 (T_2 T_1)$ proves the desired result.
\end{proof} 

\begin{exercise}{6}
Let $H$ be a Hilbert space, $T: H \to H$ is a bounded linear operator and $T^\ast$ the Hilbert-adjoint operator of $T$.
Show that $T$ is compact if and only if $T^\ast T$ is compact.
\end{exercise}
\begin{proof}
($\Rightarrow$)
Suppose $T$ is compact, so that by Lemma 8.1-2 it is bounded. 
By 8.3-2, we get that $T^\ast$ is bounded.
We can then apply Lemma 8.3-2 to $T^\ast T$ to conclude it is compact.

($\Leftarrow$)
Suppose $T^\ast T$ is compact.
Let $(x_k)$ be a bounded sequence (by $K > 0$) in $H$, so that $T^\ast T (x_n)$ contains a convergent subsequence $(T^\ast T (x_{n_k}))$.
We have 
\begin{align*}
    \norm{T x_{n_k} - T x_{n_l}}^2
    =& \brackets{T x_{n_k} - T x_{n_l}, T x_{n_k} - T x_{n_l}}\\
    =& \brackets{T (x_{n_k} - x_{n_l}), T (x_{n_k} - x_{n_l})}\\
    =& \brackets{x_{n_k} - x_{n_l}, T^\ast T (x_{n_k} - x_{n_l})}\\
    \leq& \norm{x_{n_k} - x_{n_l}} \norm{T^\ast T (x_{n_k} - x_{n_l})}\\
    \leq& 2K \norm{T^\ast T (x_{n_k} - x_{n_l})} \to 0,
\end{align*}
where the first inequality is the Schwarz inequality.
Thus, $T x_{n_k}$ is also Cauchy, and since $H$ is a Hilbert space, it converges. 
That is, $T$ is compact.
\end{proof} 

\begin{exercise}{7}
If $T$ in exercise 6 is compact, show that $T^\ast$ is compact.
\end{exercise}
\begin{proof}
We have that $T^\ast$ is bounded (3.9-2), so that $T T^\ast$ is compact (8.3-2) and, using the previous exercise (replacing $T$ by $T^\ast$ in the statement), we get that $T^\ast$ is compact, since $T T^\ast = (T^\ast)^\ast T^\ast$.
\end{proof} 

\begin{exercise}{8}
If a compact linear operator $T: X \to X$ on an infinite dimensional normed space $X$ has an inverse which is defined on all of $X$, show that the inverse cannot be bounded.
\end{exercise}
\begin{proof}
Suppose, for the sake of contradiction, that the inverse $T^{-1}$ is bounded.
Then by 8.3-2, $I = T T^{-1}: X \to X$ would be compact.
However, in 8.1-2.b, we proved that $I$ defined on an infinite dimensional space is not compact, arising a contradiction.
Thus, $T^{-1}$ must not be bounded.
\end{proof} 

\begin{exercise}{14}
Let $T: X \to X$ be a compact linear operator on a normed space.
If $\dim X = \infty$ show that $0 \in \sigma(T)$.
\end{exercise}
\begin{proof}
Suppose, for the sake of contradiction, this is not the case, so that $0 \notin \sigma(T)$.
This means that $(T - 0I)^{-1} = T^{-1}$ exists, is defined on all of $X$, and is bounded.
Thus, $T T^{-1} = I$ would be compact (8.3-2), but this contradicts 8.1-2.b, about $I$ defined on $X$ not being compact whenever $X$ is infinite dimensional.
\end{proof} 

\begin{exercise}{15}
Let $T: \ell^2 \to \ell^2$ be defined by $y = (y_i) = T x$, $x = (x_i)$, $y_{2k} = x_{2k}$, and $y_{2k - 1} = 0$, where $k = 1, 2 \dots$.
Find $\cN(T_\lambda^n)$.
Is $T$ compact?
\end{exercise}
\begin{proof}
We have that $(T - \lambda I)x = (-\lambda, x_2 - \lambda, - \lambda, x_4 - \lambda, \dots)$.
For $\lambda \neq 0$, $\cN(T_\lambda^n) = \set{0}$, for $\lambda = 0$, then $\cN(T_\lambda^n)$ would be all the vectors of the form $(x_1, 0, x_3, 0, \dots)$

Note we can think of $T$ as an operator of the form $(T x)_i = \alpha_i x_i$, where $(\alpha_n) = (0, 1, 0, 1, 0, \dots)$.
We proved (exercise 8.2.10) that such an operator would be compact if and only if $\alpha_n \to 0$, which is not the case here, so that $T$ is not compact.
\end{proof} 
